\chapter{Introducción específica} % Main chapter title

\label{Chapter2}

En este capítulo se describen los distintos protocolos de comunicación
empleados en el desarrollo del trabajo, los componentes de hardware
implementados para cumplir con los objetivos planteados y las
herramientas utilizadas para el desarrollo del software.

\section{Protocolos de comunicación}
\label{sec:protocolos_comunicacion}

Como se mencionó en la sección \ref{sec:1_4}, el dispositivo debía
cumplir con la funcionalidad de actuar como intermediario entre las
redes establecidas por Enertik y Victron, así como también constituir
un punto de acceso para que la aplicación móvil pueda operar sobre el
sistema.

Para lograr este objetivo, fue necesario implementar distintos protocolos
de comunicación que permitan el intercambio de información entre los
subsistemas involucrados, incluyendo CAN, Modbus TCP y una interfaz basada
en HTTP.

\subsection{Protocolo CAN}
El protocolo CAN (\textit{Controller Area Network}) \citep{web_intro_can}
es un protocolo serie ampliamente utilizado en la industria automotriz.
Permite la comunicación entre múltiples dispositivos sin la necesidad
de un controlador central, ofreciendo un intercambio de datos rápido
y confiable.

En este protocolo, la transmisión de datos se realiza a través de un
par de conductores denominado CAN H y CAN L, utilizando señalización
diferencial. Esta técnica consiste en transmitir la información como
la diferencia de tensión entre ambos conductores, lo que proporciona
una alta inmunidad a interferencias electromagnéticas, ya que el ruido
externo tiende a afectar a ambos cables de manera similar y, por lo
tanto, es cancelado al evaluar la diferencia entre ellos.

En cuanto a la representación lógica, un bit dominante (0 lógico) se
caracteriza por una diferencia de tensión significativa entre CAN H y
CAN L (normalmente un conductor adquiere una tensión igual a la fuente y el otro a 0 V), mientras que un bit recesivo (1 lógico) se corresponde con un estado en el cual ambos conductores presentan niveles de tensión
similares.

Por otra parte, la comunicación en una red CAN se basa en un esquema de
difusión, donde todos los nodos reciben los mensajes transmitidos. Cada
mensaje incluye un identificador que permite a los dispositivos
determinar si la información es relevante para ellos, ignorando aquellos
mensajes que no corresponden a su funcionamiento.

\subsection{Protocolo Modbus TCP}

El protocolo Modbus \citep{web_intro_modbus} es uno de los protocolos mas utilizados en la industria de la automatización. Su función principal es permitir que distintos dispositivos electrónicos se comuniquen entre sí de una forma sencilla y entandarizada.

En una red Modbus, exite una relación maestro-esclavo entre los dispositivos. El dispositivo maestro es el encargado de iniciar una comunicación y solicitar información. El esclavo se encarga de atender estas solicitudes y entregar los datos que el maestro requiera.

En el caso de Modbus TCP, este protocolo se implementa sobre redes
Ethernet, utilizando el protocolo TCP/IP como medio de transporte. Esto
permite que los dispositivos puedan comunicarse a través de redes
locales, identificándose mediante direcciones IP.

La comunicación se basa en un esquema de tipo \textit{request-response},
donde el cliente (maestro) envía una solicitud a un servidor (esclavo),
el cual procesa la petición y devuelve una respuesta. Dentro de estas
solicitudes se incluyen funciones específicas que permiten leer o
escribir distintos tipos de registros, tales como \textit{holding
registers} o \textit{input registers}.

En el contexto de este trabajo, el protocolo Modbus TCP es utilizado para
establecer la comunicación con el módulo Cerbo GX, el cual expone
diferentes registros que pueden ser consultados y
controlados de forma remota.

\subsection{Solicitudes HTTP}

El protocolo HTTP (\textit{HyperText Transfer Protocol}) \citep{web_intro_http}
es un protocolo de aplicación ampliamente utilizado para la comunicación
entre clientes y servidores. Su funcionamiento se
basa en un modelo de tipo \textit{request response}, donde un cliente
envía una solicitud (request) a un servidor, este procesa la
petición y devuelve una respuesta (response) con la información
requerida.

Las solicitudes HTTP permiten a los clientes acceder a recursos,
enviar información y ejecutar acciones sobre un servidor. Estas
peticiones están compuestas por distintos elementos, tales como el
método de la solicitud, la dirección del recurso solicitado y, en
algunos casos, un cuerpo de datos que contiene la información a enviar.

Existen diferentes tipos de métodos HTTP, cada uno con una función
específica. Entre los más utilizados se encuentran el método
GET, empleado para solicitar información al servidor, y el
método POST, utilizado para enviar datos y generar cambios en
el estado del sistema. Otros métodos relevantes incluyen
PUT, orientado a la actualización de recursos, y
DELETE, destinado a su eliminación.

En el contexto de este trabajo, las solicitudes HTTP son utilizadas
como mecanismo de comunicación entre la aplicación móvil y el
dispositivo basado en ESP32, permitiendo el envío de comandos de
control y la obtención de información del sistema.

\section{Componentes de hardware}
\label{sec:componentes_hardware}
Para el desarrollo de hardware fue necesaria la utilización de ciertos componentes. El principal de ellos fue el ESP32 como microcontrolador del equipo. A su vez, para poder adaptarlo a los distintos protocolos de comunicación, fue necesario el agregado de dos periféricos adicionales: el módulo W5110 y el MCP2551.

\subsection{Microcontrolador ESP32}
El ESP32 \citep{web_espressif} es un microcontrolador desarrollado por
la empresa Espressif Systems. Se trata de un dispositivo SoC
(\textit{System on Chip}) que integra múltiples periféricos y
capacidades de procesamiento en un único dispositivo.

Entre sus principales características, el ESP32 dispone de puertos de propósito general (GPIO), conversores analógico-digitales (ADC), módulos de comunicación serie como UART, SPI e I\textsuperscript{2}C,
temporizadores y periféricos para generación de señales PWM.

Además, el dispositivo incorpora conectividad inalámbrica, incluyendo
Wi-Fi y Bluetooth, lo que lo convierte en una plataforma adecuada para
el desarrollo de aplicaciones orientadas al campo IoT (\textit{Internet of Things}). Todas las características mencionadas anteriormente vuelven al ESP32 en una opción mas que adecuada para el desarrolo del trabajo

\subsection{Módulo transceiver}

Todos los dispositivos que pertenecen a la red de domótica de Enertik se
conectan a través de un cableado de cuatro conductores (+12 V, GND,
CAN H y CAN L). Este esquema permite tanto la alimentación como la
comunicación de los equipos conectados a la red de manera sencilla.

Por otra parte, el protocolo de comunicación empleado corresponde a una
librería propia desarrollada por Enertik, basada en el estándar CAN.
Para adaptar el ESP32 a esta red y permitir su correcta integración, fue
necesaria la implementación de un módulo transceiver.

Este módulo actúa como interfaz entre el ESP32 y la red de domótica, y
está compuesto por un microcontrolador PIC \citep{pic_datasheet}, en el
cual se implementa la librería de comunicación, y un transceptor CAN
MCP2551 \citep{2551_datasheet}, encargado de la adaptación de niveles
eléctricos para la transmisión diferencial sobre el bus.


\subsection{Módulo ethernet W5100}

Para lograr la comunicación mediante el protocolo Modbus TCP, fue
necesaria la utilización de un módulo Ethernet basado en el circuito
W5100 \citep{w5100_datasheet}.

El W5100 es un controlador Ethernet que permite dotar de conectividad
de red a sistemas embebidos, gestionando internamente la pila de
protocolos TCP/IP. La comunicación con el microcontrolador se realiza
a través de una interfaz SPI, lo que facilita su integración en
plataformas como el ESP32.

Mediante el uso de este módulo, el dispositivo desarrollado puede
establecer conexiones sobre redes Ethernet, permitiendo así la
implementación del protocolo Modbus TCP para la comunicación con
equipos del ecosistema Victron.

\section{Herramientas de software}
\label{sec:componentes_software}

Para el desarrollo del software fue necesario definir tanto el entorno
de desarrollo integrado (IDE) como el lenguaje de programación a emplear.
Asimismo, se evaluaron las herramientas necesarias para el desarrollo de
la aplicación móvil.

\subsection{IDE arduino}

El microcontrolador ESP32 es compatible con el entorno de desarrollo
Arduino, lo que permitió utilizar este IDE como plataforma principal
para el desarrollo del firmware. La elección de este entorno se basó en
su facilidad de uso, amplia documentación y la disponibilidad de una gran
cantidad de bibliotecas previamente desarrolladas.

En particular, el uso del IDE Arduino facilitó la implementación de los
distintos protocolos de comunicación requeridos en el proyecto, tales
como HTTP y Modbus TCP. Para este último, se emplearon bibliotecas
específicas que fueron adaptadas para su funcionamiento sobre el módulo
Ethernet W5100.

Adicionalmente, este entorno permite una rápida compilación y carga de
programas sobre el dispositivo, así como también herramientas de
depuración básicas, lo cual resultó adecuado para el desarrollo y prueba
del sistema.

\subsection{Funcionamiento del ESP32}

Para la programación del microcontrolador se optó por el lenguaje C++,
debido a que es el lenguaje en el cual se encuentran desarrolladas la
mayoría de las bibliotecas disponibles en el entorno Arduino, lo que
facilitó la integración de distintas funcionalidades requeridas en el
proyecto.

Como metodología de desarrollo, se utilizó FreeRTOS \citep{web_freertos},
el cual se encuentra integrado en el ESP32. Este sistema operativo en
tiempo real permite organizar el firmware mediante la ejecución de
múltiples tareas de forma concurrente, mejorando la estructura del
programa y facilitando la gestión de distintas funcionalidades en
paralelo.

En este contexto, el firmware fue dividido en distintas tareas
principales. En primer lugar, se implementó una tarea encargada de la
gestión de la comunicación Wi-Fi, donde el ESP32 opera en modo
\textit{Access Point}, permitiendo la conexión de la aplicación móvil y
actuando como servidor HTTP para la recepción de solicitudes.

Por otra parte, se desarrolló una tarea dedicada a la recepción de
mensajes provenientes de la red CAN, la cual se encarga de procesar la
información recibida desde los distintos dispositivos conectados a la
red de domótica.

Finalmente, se implementó una tarea central de procesamiento, cuya
función es interpretar la información recibida tanto desde la red CAN
como desde la interfaz HTTP, y tomar las decisiones correspondientes
para el control del sistema.

\subsection{Desarrollo de la aplicación móvil}

Para el desarrollo de la aplicación móvil se optó por utilizar el
framework Ionic \citep{web_ionic}, ya que permite la creación de aplicaciones
multiplataforma basadas en tecnologías web.

Ionic se apoya en el framework Angular, que proporciona una estructura
robusta para el desarrollo de interfaces de usuario dinámicas y
organizadas en componentes. Angular utiliza lenguajes como TypeScript,
HTML y CSS, lo que facilita la construcción de aplicaciones
modulares y escalables.

Para la adaptación de la aplicación a dispositivos móviles se utilizó
Capacitor, herramienta que permite empaquetar la aplicación web como
una aplicación nativa y facilita el acceso a funcionalidades propias
del dispositivo.

Finalmente, se empleó Android Studio para la compilación y generación
del archivo APK, permitiendo la instalación de la aplicación en
dispositivos móviles con sistema operativo Android.